\begin{helper}[Replacement gamma-boundary configuration]
Fix a defective color $\delta$, put
$\gamma=\gamma(\delta)$,
$\varphi_\delta=\operatorname{swap}_{\gamma,\delta}
 (\varphi;v^0_\delta)$, let
$X_\delta=V(T_\delta)$, and put
$Z=V(T)\setminus U_\gamma$.  The color $\gamma$ is defective for
$X_\delta$ under $\varphi_\delta$.  The three pairwise distinct edge
identities
\[
 f^0_\delta,qquad f^1_\gamma,qquad f^2_\gamma
\]
all have color $\gamma$ under $\varphi_\delta$, and they are exactly the
$\gamma$-colored edges with one end in $Z$ and the other outside
$X_\delta$.
\end{helper}
\begin{proof}
By \noderef{ScheideDefectiveSwitchEffect}, the switched coloring gives
$f^0_\delta$ color $\gamma$ and fixes every other identity incident with
$V(T)$.  The other end of $f^0_\delta$ is outside $V(T)$ by
\noderef{ScheideDefectiveSetup}, whereas each displayed edge of the
deleted $\gamma$-block has both ends in $V(T)$ by
\noderef{ScheideStructuralConfiguration}; hence both block identities
differ from $f^0_\delta$ and retain their structural color $\gamma$.
The two block identities are distinct by exact-edge injectivity in
\noderef{TreeSequence}.  Their base ends lie in the
base $Y\subseteq Z$ and their block ends lie in
$U_\gamma=V(T)\setminus X_\delta$, using the exact vertex-set equality in
\noderef{ScheideReplacementData}.

The edge $f^0_\delta$ also leaves $X_\delta$.  Its tree end
$v^0_\delta$ is an exit by \noderef{ScheideDefectiveSetup}, hence belongs
to $Y\subseteq Z$ by \noderef{ScheideStructuralExitsInBase}.  Its other
end cannot equal $u^1_\delta$ or $u^2_\delta$: otherwise
$f^0_\delta$ and the corresponding distinct edge $f^i_\delta$ would be
two incident edges of original color $\delta$, contrary to the properness
in the globally maximum triple.  The replacement vertex-set equality now
puts that other end outside $X_\delta$.

Thus the three identities are pairwise distinct boundary edges of
$X_\delta$ and already show, by \noderef{DefectiveColors}, that
$\gamma$ is defective.

It remains to prove exactness.  Let $a$ be a $\gamma$-colored edge under
$\varphi_\delta$ with one end in $Z$ and the other outside $X_\delta$.
If $a=f^0_\delta$ there is nothing to prove.  Otherwise the pointwise
fixation clause of \noderef{ScheideDefectiveSwitchEffect} gives
$\varphi(a)=\gamma$.  Original closedness from
\noderef{ScheideGlobalMaximumOrder}, together with the fact that
$\gamma$ is missing at $v^0_\delta$, excludes an original
$\gamma$-edge from $Z$ to the exterior of $V(T)$.  Its other endpoint
therefore lies in $V(T)\setminus X_\delta=U_\gamma$.  At either vertex
of this pair, the corresponding displayed block edge is already incident
and has color $\gamma$; original properness makes it the unique such
identity.  Hence $a=f^1_\gamma$ or $a=f^2_\gamma$.  Conversely the
endpoint and color calculations above show directly that each of the
three named identities has the stated departure property.
\end{proof}
